\documentclass{article}
\usepackage[utf8]{inputenc}
\usepackage{geometry}
\geometry{a4paper, margin=1in}
\usepackage{hyperref}
\usepackage{listings}
\usepackage{xcolor}

\definecolor{codegreen}{rgb}{0,0.6,0}
\definecolor{codegray}{rgb}{0.5,0.5,0.5}
\definecolor{codepurple}{rgb}{0.58,0,0.82}
\definecolor{backcolour}{rgb}{0.95,0.95,0.92}

\lstdefinestyle{mystyle}{
    backgroundcolor=\color{backcolour},   
    commentstyle=\color{codegreen},
    keywordstyle=\color{magenta},
    numberstyle=\tiny\color{codegray},
    stringstyle=\color{codepurple},
    basicstyle=\ttfamily\footnotesize,
    breakatwhitespace=false,         
    breaklines=true,                 
    captionpos=b,                    
    keepspaces=true,                 
    numbers=left,                    
    numbersep=5pt,                  
    showspaces=false,                
    showstringspaces=false,
    showtabs=false,                  
    tabsize=2
}
\lstset{style=mystyle}

\title{Technical Execution Plan: Online Scene Change Detection (O-SCD)\\
\large Custom Dataset Generation and Training Pipeline}
\author{Joshua Hecke}
\date{\today}

\begin{document}

\maketitle

\section{Introduction}
This document outlines the end-to-end technical steps for creating a custom time-variant dataset and training the \textbf{O-SCD} (Online Scene Change Detection) framework. Because we do not yet have a multi-temporal underwater dataset, we must construct a custom dataset consisting of two states: a reference scene (Object A) and an inference/query scene (Object A + Object B).

O-SCD operates in a label-free, online manner by relying on a pre-computed 3D Gaussian Splatting (3DGS) reconstruction of the \textit{reference scene}, and performs PnP-based pose estimation on the fly for the \textit{inference scene}.

\section{Phase 1: Dataset Generation and Structure}
To successfully run O-SCD, your custom dataset must adhere exactly to the directory structure expected by the dataloaders.

\subsection{Image Capture}
\begin{enumerate}
    \item \textbf{Reference Scene (Time 1):} Capture a set of images (e.g., 50-150) covering the target object/area (Object A). Save these images into a folder named \texttt{input} inside a \texttt{reference\_scene} directory.
    \item \textbf{Inference Scene (Time 2):} Introduce the change (Object B) into the same environment. Capture a new set of images from varying perspectives. Save these images directly into an \texttt{images} folder inside an \texttt{inference\_scene} directory.
\end{enumerate}

\subsection{Expected Directory Layout}
Before running any scripts, organize your custom dataset as follows:
\begin{lstlisting}[language=bash]
~/UMP/Codebase/Dataset/Custom_OSCD_Dataset/
|--- reference_scene/
|    |--- input/
|         |--- frame001.jpg
|         |--- frame002.jpg ...
|--- inference_scene/
     |--- input/
          |--- frame001.jpg
          |--- frame002.jpg ...
\end{lstlisting}

\section{Phase 2: Preparation and COLMAP Processing}
\textbf{CRITICAL PREPROCESSING STEP:} O-SCD requires camera poses \textbf{only} for the reference scene (poses for the inference scene are estimated on the fly via PnP). However, because O-SCD directly inherits the \textit{undistorted} pinhole camera intrinsics from the reference scene and applies them to the inference scene, \textbf{the inference images must also be undistorted} to prevent systematic projection errors. 

The easiest way to achieve this is to run the standard 3DGS \texttt{convert.py} on \textbf{both} scenes independently. This will generate undistorted images in an \texttt{images/} folder for both.

\begin{lstlisting}[language=bash]
# Use the standard 3DGS conversion tool to run COLMAP
cd ~/UMP/Codebase/Tools/gaussian-splatting-main
conda activate colmap_runner

# Generate sparse point cloud and undistorted images for the reference scene
python convert.py -s ~/UMP/Codebase/Dataset/Custom_OSCD_Dataset/reference_scene

# Generate undistorted images for the inference scene (O-SCD will ignore its sparse poses)
python convert.py -s ~/UMP/Codebase/Dataset/Custom_OSCD_Dataset/inference_scene
\end{lstlisting}
This will create \texttt{sparse/0} and \texttt{images/} inside both directories. O-SCD's dataloader is hardcoded to load from \texttt{inference\_scene/images}, which will now correctly contain undistorted images.

\section{Phase 3: Reference Scene Reconstruction}
O-SCD requires a fully trained 3DGS point cloud of the reference scene as a prior. It expects this reconstruction to be stored in a directory named \texttt{reference\_reconstruction} at the root of the dataset.

\begin{lstlisting}[language=bash]
# Switch to the standard 3DGS environment
conda activate 3dgs # (or whichever env you use for standard 3DGS)

# Train the baseline model on the reference scene
# By setting -m to reference_reconstruction, it saves directly where O-SCD expects it
python train.py -s ~/UMP/Codebase/Dataset/Custom_OSCD_Dataset/reference_scene \
                -m ~/UMP/Codebase/Dataset/Custom_OSCD_Dataset/reference_reconstruction
\end{lstlisting}
Verify that \texttt{\textasciitilde/UMP/Codebase/Dataset/Custom\_OSCD\_Dataset/reference\_reconstruction/point\_cloud/iteration\_30000/point\_cloud.ply} exists.

\section{Phase 4: Running O-SCD}
With the environment prepared, you can now run the O-SCD pipeline.

\subsection{Environment Setup}
\begin{lstlisting}[language=bash]
cd ~/UMP/Codebase/3DGS-Change-Detection/O-SCD-main

# Initialize and update submodules (CRITICAL for requirements.txt to succeed)
git submodule update --init --recursive

conda create -n oscd python=3.12 -y
conda activate oscd

# Install dependencies (CUDA 12.x recommended based on RTX 5070 Ti)
pip install torch torchvision xformers --index-url https://download.pytorch.org/whl/cu121
pip install cupy-cuda12x
pip install -r requirements.txt
\end{lstlisting}

\subsection{Executing Online Change Detection}
Run the core \texttt{oscd.py} script. This script loads the reference point cloud, aligns the inference images, detects the changes, and generates the change masks.

\begin{lstlisting}[language=bash]
# Run Change Detection
# --test_hold 5 holds out 1/5th of the frames for testing (if needed)
# --resolution 1 ensures full resolution (default scripts use 4 for downsampling)
python oscd.py -s ~/UMP/Codebase/Dataset/Custom_OSCD_Dataset \
               -m ~/UMP/Codebase/Dataset/Custom_OSCD_Dataset/output \
               --resolution 1 \
               --test_hold 5
\end{lstlisting}

\subsection{Executing Scene Update (Optional)}
O-SCD also provides an update script to modify the 3D Gaussian Splatting scene representation based on the detected changes.
\begin{lstlisting}[language=bash]
python update.py -s ~/UMP/Codebase/Dataset/Custom_OSCD_Dataset \
                 -m ~/UMP/Codebase/Dataset/Custom_OSCD_Dataset/output \
                 --resolution 1
\end{lstlisting}

\end{document}
